% =====================================================================
%  1bat-ud3-1.tex  —  HORA 1: Què és una funció? Relacions de la vida quotidiana
%  (sense preàmbul: s'inclou des de main.tex)
% =====================================================================
\horatitol{Sessió 1 — Què és una funció? Relacions de la vida quotidiana}%
{De les relacions de dependència quotidianes a la idea de funció: una magnitud que depèn d'una altra.}

\subsection{Idea clau}

A la UD2 vam aprendre a posar una \textbf{lletra} a una quantitat desconeguda i a
escriure expressions com $2x$ o $8x+100$. Ara fem un pas més: moltes vegades una
magnitud \textbf{depèn} d'una altra. El cost de la llum depèn dels kWh consumits;
l'ajut que rep cada família depèn de quantes famílies se'l reparteixen. Quan una
magnitud queda determinada per una altra, diem que hi ha una \textbf{funció}.

\begin{idea}
\textbf{Què és una funció (de manera intuïtiva)}\\
Una \textbf{funció} és una relació entre dues magnituds en què, per a cada valor de
la primera (la variable de què \emph{depenem}, sovint $x$), n'obtenim
\textbf{un de determinat} de la segona (la que \emph{depèn}, sovint $y$).
\begin{itemize}[leftmargin=1.4em, itemsep=2pt]
  \item «$y$ \textbf{depèn de} $x$» \quad s'escriu \quad $y=f(x)$.
  \item La mateixa dependència es pot llegir de \textbf{tres maneres}: amb una
        \textbf{frase}, amb una \textbf{taula} de valors o amb una \textbf{gràfica}.
\end{itemize}
\end{idea}

\subsection{Exemples}

\begin{exemple}[el cost del menjador depèn dels àpats]
Un menjador comunitari cobra $2\eur$ per àpat. La recaptació $R$ \textbf{depèn} del
nombre d'àpats $x$:
\[
R = f(x) = 2x \quad(\text{en euros}).
\]
Per a cada valor de $x$ n'obtenim \emph{un} de $R$: si $x=10$, aleshores $R=20\eur$;
si $x=120$, aleshores $R=240\eur$. La taula de valors ho deixa clar:

\begin{center}
\begin{tabular}{c c c c}
\toprule
Àpats $x$ & $10$ & $50$ & $120$ \\
\midrule
Recaptació $R$ ($\eur$) & $20$ & $100$ & $240$ \\
\bottomrule
\end{tabular}
\end{center}
Aquí $x$ és la \textbf{variable independent} (la que decidim) i $R$ la
\textbf{variable dependent} (la que en surt).
\end{exemple}

\begin{exemple}[l'ajut per família depèn del nombre de famílies]
Un fons d'emergència de $\num{1200}\eur$ es reparteix a parts iguals entre les
famílies afectades. Els diners $D$ que rep cada família \textbf{depenen} del nombre
de famílies $x$:
\[
D = f(x) = \frac{\num{1200}}{x} \quad(\text{en euros}).
\]
Si $x=2$ famílies, cadascuna rep $600\eur$; si $x=4$, en rep $300\eur$; si n'hi ha
moltes, cadascuna en rep cada cop menys. Una mateixa relació, escrita en una sola
expressió, val per a \emph{qualsevol} nombre de famílies.
\end{exemple}

\subsection{Exercicis resolts}

\begin{exresolt}[identificar qui depèn de qui]
En cada situació, digues quina magnitud és la \textbf{variable independent} ($x$) i
quina és la \textbf{dependent} ($y$):
\begin{enumerate}[label=\alph*), leftmargin=1.6em]
  \item El cost mensual de la llum segons els kWh consumits.
  \item El temps total d'atenció segons el nombre de visites a domicili.
  \item Els diners recaptats en una rifa solidària segons els bitllets venuts.
\end{enumerate}
\solucio
La variable independent és \emph{la causa} (allò que fem variar) i la dependent és
\emph{el resultat} (allò que en surt):
\begin{enumerate}[label=\alph*), leftmargin=1.6em]
  \item Independent: kWh consumits. Dependent: cost de la llum.
  \item Independent: nombre de visites. Dependent: temps total d'atenció.
  \item Independent: bitllets venuts. Dependent: diners recaptats.
\end{enumerate}
\resposta{En els tres casos, el \emph{resultat} (cost, temps, diners) depèn de la \emph{causa} (kWh, visites, bitllets).}
\end{exresolt}

\begin{exresolt}[de la frase a la taula]
Una treballadora familiar fa visites a domicili de $45$ minuts cadascuna. El temps
total $T$ depèn del nombre de visites $x$. Escriu la relació com a funció i completa
la taula per a $x=1,2,3,4$ visites.
\solucio
Cada visita són $45$ minuts, i en fa $x$, per tant $T=f(x)=45x$ (en minuts).
Substituïm cada valor de $x$:
\[
f(1)=45,\quad f(2)=90,\quad f(3)=135,\quad f(4)=180.
\]
\begin{center}
\begin{tabular}{c c c c c}
\toprule
Visites $x$ & $1$ & $2$ & $3$ & $4$ \\
\midrule
Temps $T$ (min) & $45$ & $90$ & $135$ & $180$ \\
\bottomrule
\end{tabular}
\end{center}
\resposta{$T=f(x)=45x$; la taula dóna $45,\ 90,\ 135,\ 180$ minuts.}
\end{exresolt}

\subsection{Exercicis per fer}

\begin{exercicis}
\begin{exlist}
  \item En cada cas, indica quina és la variable \textbf{independent} i quina la
        \textbf{dependent}:
    \begin{enumerate}[label=\alph*), leftmargin=1.6em, topsep=2pt]
      \item El preu total de la compra segons el nombre de productes.
      \item La nota d'un treball segons les hores dedicades.
      \item L'aigua que consumeix un casal segons el nombre d'infants.
    \end{enumerate}
  \item Un casal cobra $4\eur$ per infant i activitat. El cost total $C$ depèn del
        nombre d'infants $x$. Escriu $C=f(x)$ i calcula $f(5)$, $f(20)$ i $f(35)$.
  \item Una associació ven samarretes solidàries a $8\eur$ cadascuna. Els diners
        recaptats $D$ depenen de les samarretes venudes $x$. Escriu $D=f(x)$ i
        completa una taula per a $x=10,25,50$.
  \item Un fons de $\num{600}\eur$ es reparteix a parts iguals entre $x$ persones
        voluntàries. Els diners $D$ que rep cadascuna depenen de $x$. Escriu
        $D=f(x)$ i calcula què rep cada persona si $x=3$, $x=6$ i $x=12$.
  \item \textbf{(Interpretació)} Mira la taula següent, que relaciona el nombre
        d'usuaris d'un servei amb el cost mensual:
        \begin{center}
        \begin{tabular}{c c c c}
        \toprule
        Usuaris $x$ & $10$ & $20$ & $30$ \\
        \midrule
        Cost $C$ ($\eur$) & $250$ & $450$ & $650$ \\
        \bottomrule
        \end{tabular}
        \end{center}
        Quants euros augmenta el cost cada cop que entren $10$ usuaris nous? Quin
        creus que seria el cost amb $0$ usuaris?
  \item \textbf{(Raonament)} Escriu \emph{amb les teves paraules} què és per a tu una
        funció i posa un exemple propi, de l'àmbit social o de la vida diària, de
        dues magnituds en què una depengui de l'altra. Indica quina és la variable
        independent i quina la dependent.
\end{exlist}
\end{exercicis}

% ---------------------------------------------------------------------
\fullsolucions{Sessió 1}
\begin{solucions}
\begin{sollist}
  \item (a) Indep.: nombre de productes; dep.: preu total. \quad
        (b) Indep.: hores dedicades; dep.: nota. \quad
        (c) Indep.: nombre d'infants; dep.: aigua consumida.
  \item $C=f(x)=4x$; \quad $f(5)=20\eur$, \quad $f(20)=80\eur$, \quad $f(35)=140\eur$.
  \item $D=f(x)=8x$; \quad taula: $x=10\rightarrow 80\eur$, \quad
        $x=25\rightarrow 200\eur$, \quad $x=50\rightarrow 400\eur$.
  \item $D=f(x)=\dfrac{\num{600}}{x}$; \quad $f(3)=200\eur$, \quad
        $f(6)=100\eur$, \quad $f(12)=50\eur$.
  \item El cost augmenta $200\eur$ cada $10$ usuaris nous. Amb $0$ usuaris, el cost
        seria $50\eur$ (la part fixa: $250-200=50$).
  \item Resposta oberta. Exemple vàlid: «els diners d'una guardiola segons les
        setmanes que hi estalvio» $\rightarrow$ indep.: setmanes; dep.: diners.
\end{sollist}
\end{solucions}
